\chapter{Decision register}
\label{ch:decisions}

\begin{quote}\itshape
Every material free choice made in this study, its value, and the \emph{kind} of justification
behind it. Justification classes: \textbf{[M]}~measured/certified data; \textbf{[L]}~literature
range; \textbf{[D]}~derived from our own model; \textbf{[C]}~community convention or benchmark
practice; \textbf{[T]}~engineering default, explicitly exposed to the symmetric tuning budget
of phase F5; \textbf{[A]}~assumption with declared sensitivity. A choice justified by none of
these classes would be arbitrary --- the audit that produced this table found two such cases,
both already resolved (noted below).
\end{quote}

\section*{What this register is, and how to use it}

Every quantitative study rests on choices that the data did not make for it: which engine
variant to model, how much noise to inject, where to place a detection threshold, how many
tuning trials to allow each method. Most papers leave these choices implicit, scattered
through the text, or simply unstated --- and that is precisely what makes replication hard.
This appendix collects \emph{every material free choice} in the study in one place, states
the value chosen, and --- most importantly --- classifies \emph{why} that value is defensible.

The classification is the point. A value is not more or less true because it appears in a
table; what matters is the \emph{kind} of evidence behind it, because the kind tells a reader
what would need to change for the choice to change:

\begin{description}
  \item[{[M]} Measured or certified.] The value comes from an authoritative measurement ---
    a type-certificate data sheet, the ICAO emissions databank. Example: the calibration
    anchors of \cref{tab:anchors}. To challenge it, you would challenge the source document.
  \item[{[L]} Literature range.] The value sits inside a range published by independent
    studies. Example: fouling magnitudes from the gas-turbine degradation literature
    (\cref{tab:degradation}). Challenging it means producing better field data.
  \item[{[D]} Derived.] Computed from our own calibrated model, so it inherits the model's
    validation. Example: the rating temperature $T_{4,\max}$, fixed by the takeoff anchor.
  \item[{[C]} Convention.] A community default adopted for comparability, not because it is
    uniquely correct. Example: the M0.78/FL350 design point. Deviating is legitimate but
    costs comparability with other studies.
  \item[{[T]} Tunable default.] An engineering guess that we ourselves did not trust ---
    so it was explicitly exposed to the symmetric tuning budget of phase F5, where both
    families could revise it. Example: detector window lengths. These are the choices the
    study \emph{stress-tested} rather than defended.
  \item[{[A]} Assumption.] A value we could not ground in any of the above; declared openly,
    with its sensitivity noted. These are the study's honest soft spots.
\end{description}

To use the register for replication: reproduce the [M] and [L] values from their cited
sources, accept the [D] values if you accept the model calibration of \cref{ch:model},
adopt or revise the [C] conventions to taste, re-tune the [T] defaults under your own budget,
and probe the [A] assumptions first if you suspect a result. A reader who disagrees with a
conclusion should be able to find, in this table, the exact choice they would have made
differently --- and the class tells them what evidence their alternative needs.

\section*{Engine model (F1)}

\begin{center}\footnotesize\setlength{\tabcolsep}{4pt}
\begin{tabular}{>{\raggedright\arraybackslash}p{0.30\textwidth}>{\raggedright\arraybackslash}p{0.15\textwidth}c>{\raggedright\arraybackslash}p{0.38\textwidth}}
  \toprule
  Decision & Value & Class & Basis \\
  \midrule
  Engine variant & CFM56-7B26 & [C] & most common -7B; richest public data \\
  Design point & M0.78 / FL350 / ISA & [C] & aerodynamic-design-point convention \\
  Calibration targets \& tolerances & \cref{tab:tcds}; $\pm3$/$\pm5$\,\% & [M] & TCDS + EEDB; tolerances set a priori in the spec \\
  Cruise thrust/TSFC targets & 5480 lbf / 0.627 & [A] & public type data, still flagged
    TO\_VERIFY in the contract; H1 was gated on measured anchors only \\
  HPC design PR & 9.35 & [M] & calibrated to the measured SLS OPR 27.61 \\
  HP map speed reference & 13 940 rpm & [M] & placed so rated-thrust N2 respects the TCDS redline \\
  Rating temperature $T_{4,\max}$ & 3200\,°R & [D] & from the calibrated takeoff anchor (\cref{sec:decks}) \\
  ICM perturbation step & $\pm0.5$\,\% & [C]/[D] & central-difference standard: large enough to
    dominate the $10^{-8}$ solver tolerance, small enough for linearity --- and, since
    L-ICM, \emph{measured} rather than argued: swept against $\pm0.25$/$\pm1.0$\,\%
    (\cref{sec:f-icm}) \\
  Confusable threshold & $15^{\circ}$ & [T] & pre-registered definition for H2; sensitivity
    checked implicitly by reporting all angles \\
  Deck envelope blocks \& PC $\geq 0.5$ & \cref{sec:decks} & [C] & EHM trending conditions \\
  Thermo method & tabular & [D] & CEA cross-check within 0.7\,\% (\cref{tab:cea}) \\
  \bottomrule
\end{tabular}
\end{center}

\section*{Synthetic fleet (F2)}

\begin{center}\footnotesize\setlength{\tabcolsep}{4pt}
\begin{tabular}{>{\raggedright\arraybackslash}p{0.30\textwidth}>{\raggedright\arraybackslash}p{0.15\textwidth}c>{\raggedright\arraybackslash}p{0.38\textwidth}}
  \toprule
  Decision & Value & Class & Basis \\
  \midrule
  Degradation magnitudes/profiles & \cref{tab:degradation} & [L] & Diakunchak; Kurz \& Brun \\
  Fouling time constant & 3000 cycles & [L] & fouling saturates within months of service \\
  Wash interval / recovery & 800--1600 cy / 30--70\,\% & [L] & typical wash programs;
    recovery range from the same sources \\
  New-engine EGT margin & $85\pm6$\,°C & [L] & public -7B two-digit margins; per-engine
    scatter is a manufacturing-variability assumption [A] \\
  Severity multiplier spread & LogNormal $(0,\,0.3)$ & [A] & produces the observed-in-service
    spread of lives (\cref{tab:fleet}); a fleet-realism knob, not a fitted quantity \\
  Sensor noise / quantization & catalog table & [L] & typical ACARS-class instrumentation \\
  Acute-fault targets \& magnitudes & 6 params, 0.4--1.5\,\% & [D] & drawn from the ICM
    confusable set; magnitudes bracket the detectability edge (sub-noise to clearly visible) \\
  Acute probability / onset & 0.5 / 30--90\,\% of life & [T] & enough episodes per class for
    test statistics; onset as life fraction after the audit-found bug (\cref{sec:audit-difficulty}) \\
  Ambient scatter (takeoff dTs) & $\mathcal{N}(10, 8)$\,°C clipped & [A] & plausible mixed
    fleet climate; affects realism, not comparison fairness (common to both pipelines) \\
  Difficulty gate threshold & 60\,\% & [T] & pre-set before the audit ran; its bite was
    demonstrated when the first design failed it \\
  Splits & 70/10/20 by engine & [C] & leakage-free standard; CI-checked \\
  \bottomrule
\end{tabular}
\end{center}

\section*{Traditional pipeline (F3) and evaluation}

\begin{center}\footnotesize\setlength{\tabcolsep}{4pt}
\begin{tabular}{>{\raggedright\arraybackslash}p{0.30\textwidth}>{\raggedright\arraybackslash}p{0.15\textwidth}c>{\raggedright\arraybackslash}p{0.38\textwidth}}
  \toprule
  Decision & Value & Class & Basis \\
  \midrule
  Holt gains & $\alpha{=}0.15, \beta{=}0.05$ & [T] & responsive-but-smooth default; swept in F5 \\
  CUSUM parameters & $k{=}0.75\sigma, h{=}8$ & [T] & standard allowances; swept in F5 \\
  Gap-detector EWMAs & 0.1 / 0.005 & [T] & fast follows a ramp, slow spans a wash cycle; swept in F5 \\
  Alarm persistence & 10-of-14 & [T] & multi-sample confirmation; swept in F5 \\
  Kalman $q$, prior, $\lambda$ & $2{\cdot}10^{-4}$, 2\,\%, 1 & [T] & $q$ sized to chronic
    rates; all swept in F5 \\
  Wash reset fraction & 0.5 & [A] & mid-range of the catalog's 30--70\,\% recovery \\
  RUL window / cap & 1500 cy / 25 000 cy & [T]/[C] & window spans a wash cycle; cap is
    NASA-benchmark practice \\
  NASA-score scaling & $d/100$ cycles & [A] & the PHM08 constants assume C-MAPSS-scale RUL;
    scaling declared, identical for both families \\
  Assumed nominal margin & 85\,°C & [A] & fleet nominal; per-engine estimation is queued as a
    pipeline improvement in F5 \\
  Statistical tests & Wilcoxon, Holm, BCa & [C] & pre-registered (\cref{sec:safeguards}) \\
  RUL cap (AI) & 12\,000 cy & [C] & piecewise-linear RUL target, standard prognostics practice; same cap both families in F5 \\
  GRU/AE architectures \& gains & \cref{ch:ai} & [T] & engineering defaults; swept in F5 \\
  Gradient-boosting impl. & sklearn HistGB & [D] & XGBoost's second OpenMP runtime segfaults beside torch-MPS on macOS; same model family \\
  Conformal level & 90\,\% & [C] & matches hypothesis H5's coverage target \\
  PCS attribution basis & SHAP of predicted class, step block, cockpit channels & [T] & one defensible instantiation of C5; alternatives (true-class attribution, all blocks) belong to the F5 sensitivity study \\
  Detector sweep budget & 20 configs, $\leq$1 val false alarm & [T] & pre-F5 threshold exercise; selection on validation only \\
  Sim-to-real benchmark & C-MAPSS FD001 & [C]/[A] & plan named N-CMAPSS; multi-GB distribution breaks the desktop norm --- substitution to the canonical 5-MB benchmark, disclosed; not frozen by prereg-v1 \\
  FD001 RUL cap & 130 cycles & [C] & community standard for that benchmark \\
  $\alpha$--$\lambda$ settings & $\alpha{=}20\,\%$, $\lambda{=}\{.5,.7,.9\}$ & [C] & common
    prognostics practice \\
  \bottomrule
\end{tabular}
\end{center}

\section*{Later phases (F5--F10) and economics}

\begin{center}\footnotesize\setlength{\tabcolsep}{4pt}
\begin{tabular}{>{\raggedright\arraybackslash}p{0.26\textwidth}>{\raggedright\arraybackslash}p{0.17\textwidth}c>{\raggedright\arraybackslash}p{0.40\textwidth}}
  \toprule
  Decision & Value & Class & Basis \\
  \midrule
Tuning budget & 50 Optuna (TPE) trials per family & [C] & symmetric by design; enough for a
    14--16-dim space, disclosed in prereg-v1 \\
  Detection selection rule & lexicographic min-FA, then max-recall & [D] & no trial met the
    FA$\leq$1 val budget; infeasibility rule disclosed in prereg-v1 \\
  H1 operating point & per-family, FA-budgeted & [D] & fixed-recall-0.9 unreachable at v1.1
    magnitudes; amended \emph{at} freeze, on the record \\
  Report-schedule design (F7) & greedy over the ICM grid & [M] & the schedule is an
    observability design variable; measured, not assumed \\
  Learned-MOPA architecture & GRU over WLS projections & [D] & raw-sequence GRU
    failed; physics-inversion-first is the working split \\
  H7.3 drift thresholds & $\geq$20 / $\leq$10 pp & [T] & calibrated on a 16-episode dev sample;
    refuted as frozen --- discipline over narrative \\
  Surrogate architecture (L1) & linearization $+$ MLP residual & [D] & pure MLP could not match
    N2 precision; residual form did \\
  Surrogate gate & like-for-like on the same test set & [D] & two earlier drafts (statistic
    mismatch, wrong population) disclosed in code \\
  Fleet-v2 emitter & surrogate (flag) & [D] & v1.1 left byte-identical and frozen;
    v2 in a separate directory \\
  Twin-residual feature (L6) & $(\mathbf{I}-\mathbf{H}_{\text{ref}}\mathbf{M})\Delta\mathbf{z}$
    & [M] & the residual the linear GPA cannot explain; H4 refuted a 2nd way \\
  M3 (physics loss) & deferred & [D] & two H4 negatives shift the burden; needs its own
    pre-registration \\
  ICM robustness to calibration & measured, not assumed & [D] & the premise that the study
    rests on the ICM's \emph{structure}, not the twin's absolute values, is tested by
    perturbing every calibration knob inside its contract tolerance (L-ICM, prereg-v13,
    \cref{sec:f-icm}) \\
  Certificate instrument (F10) & first-order Fisher (linear $\mathbf{H}$) & [D] & curvature
    below the noise floor (measured), so first-order governs identifiability \\
  Certificate region scale & CRB shape $+$ split-conformal scale & [D] & CRB is a lower bound
    for unbiased estimators; conformal supplies the absolute scale \\
  Economics: single mechanism & unsched.\ $\to$ sched.\ & [D] & the only
    mechanism the results directly support; all else discussed, not monetized \\
  RUL method (traditional) & Theil--Sen linear extrapolation & [C]/[A] & operational-practice baseline (what airlines field), not research SOTA (particle filters, similarity-based); a fair industry baseline, disclosed; an advanced method would narrow but not close H3 (F8/L-RUL) \\
    Economics: absolute anchor & industry unscheduled band, not the raw model & [A] & the raw
    synthetic-fleet rate overstates practice; disclosed \\
\bottomrule
\end{tabular}
\end{center}

\paragraph{Audit outcome.} Sweeping the study for unjustified choices found two genuinely
arbitrary decisions, both caught and resolved before this appendix was written: the
\emph{chronic-mechanism labeling} of the diagnosis task (resolved by the acute-episode
redesign, \cref{sec:audit-difficulty}) and the \emph{acute onset drawn against the simulation
horizon} rather than engine life (resolved by two-pass generation). Remaining class-[A]
assumptions are declared at their point of use and shared by both EHM families, so they affect
the realism of the world, not the fairness of the comparison. Class-[T] defaults are the
traditional pipeline's tuning surface --- left untouched by design until F5 spends the same
optimization budget on both families. The later-phase rows (F5--F10 and economics) were added
as those phases ran; each material choice there is either measured [M], a disclosed
convention [C], a reasoned design decision [D] traceable to a result, or a flagged assumption
[A] shared by both families or confined to the economic model. Two [T] thresholds were
themselves refuted as frozen (H7.3, and H10.2's $\chi^2$ radius), which the register treats
as evidence the gates bite rather than as choices to relitigate.
